\section*{Câu 4.4.1 --- K-Nearest cho quyết định cho vay}

Đề cho bảng tám khách huấn luyện (Tuổi, Gia đình, Thu nhập, Số tiền vay, Kết quả \textbf{Yes}/\textbf{No}) và một khách mới có thông tin: Tuổi $23$, Gia đình \textbf{Có}, Thu nhập $32$ triệu, vay $250$ triệu. Ta dùng thuật toán láng giềng gần nhất và \emph{tự chọn} công thức khoảng cách cùng giá trị $K$. Ở đây ta chọn $K=5$ và khoảng cách Euclidean trên vector bốn chiều, sau khi mã hoá ``Gia đình'' thành số.

\paragraph{Mã hoá và vector}
Quy ước \textbf{Có}$\,=1$, \textbf{Không}$\,=0$. Lượng thu nhập và số tiền vay giữ nguyên đơn vị như trong bảng (triệu). Với một mẫu có vector đặc trưng $(x_1,x_2,x_3,x_4)$ với $x_2\in\{0,1\}$:
\[
\mathbf{x}=\bigl(\text{Tuổi},\, x_2,\,\text{Thu nhập},\,\text{Số tiền vay}\bigr).
\]
Điểm cần dự báo có vector
\[
\mathbf{q}=\bigl(23,\,1,\,32,\,250\bigr).
\]
Khoảng cách giữa $\mathbf{x}$ và $\mathbf{q}$ là
\[
  d(\mathbf{x},\mathbf{q})=\sqrt{\sum_{\ell=1}^4(x_\ell-q_\ell)^2},
  \qquad
  d^{2}(\mathbf{x},\mathbf{q})=\sum_{\ell=1}^4(x_\ell-q_\ell)^{2}.
\]
\emph{(Ở dưới ta ghi chủ yếu $d^2$ cho gọn --- cột $d=\sqrt{d^{2}}$ chỉ để so sánh độ gần.).}

\begin{center}
{\small
\begin{tabular}{@{}clr@{}}\toprule
ID & Vector $(t_1,t_2,t_3,t_4)$ --- Tuổi, Gia đình (0/1), Thu nhập, Vay & Kết quả \\\midrule
1 & $(25,\,1,\,20,\,120)$ & Yes\\
2 & $(22,\,0,\,15,\,200)$ & No\\
3 & $(24,\,1,\,18,\,180)$ & Yes\\
4 & $(30,\,0,\,17,\,130)$ & Yes\\
5 & $(31,\,1,\,22,\,250)$ & No\\
6 & $(35,\,1,\,22,\,150)$ & Yes\\
7 & $(22,\,0,\,23,\,210)$ & No\\
8 & $(28,\,1,\,19,\,280)$ & No\\\bottomrule
\end{tabular}
}
\end{center}

\paragraph{Ví dụ một phép tính tay (ID 5 --- gần nhất)}
Độ lệch theo các chiều so với $\mathbf{q}$:
\[
(23-31,\,1-1,\,32-22,\,250-250)=(-8,0,10,0).
\]
Do đó $d^{2}_{5}=8^{2}+0^{2}+10^{2}+0^{2}=64+100=164$, $d_{5}=\sqrt{164}\approx 12{,}81$.

\paragraph{Bình phương khoảng cách tới các mẫu còn lại và sắp xếp tăng dần theo khoảng cách.}
Tính tương tự được bảng (đã làm tròn $d=\sqrt{d^2}$):

\begin{center}
{\small
\begin{tabular}{@{}crccl@{}}\toprule
TT & ID & $d^{2}=d(\cdot,\mathbf{q})^{2}$ & $d$ & Nhãn \\\midrule
1 & 5 & $164$ & $\approx 12{,}81$ & No\\
2 & 8 & $1094$ & $\approx 33{,}08$ & No\\
3 & 7 & $1683$ & $\approx 41{,}02$ & No\\
4 & 2 & $2791$ & $\approx 52{,}83$ & No\\
5 & 3 & $5097$ & $\approx 71{,}39$ & Yes\\
6 & 6 & $10244$ & $\approx 101{,}21$ & Yes\\
7 & 4 & $14675$ & $\approx 121{,}14$ & Yes\\
8 & 1 & $17048$ & $\approx 130{,}57$ & Yes\\\bottomrule
\end{tabular}
}
\end{center}

\paragraph{Chọn $K$ và bỏ phiếu đa số}
Ta chọn $K=5$ (số lẻ, thường dùng trong k-NN). Năm láng giềng gần nhất theo bảng trên là ID $5,8,7,2,3$ với nhãn lần lượt
\[
\text{No},\ \text{No},\ \text{No},\ \text{No},\ \text{Yes}.
\]
Đếm phiếu: \textbf{No} có bốn phiếu, \textbf{Yes} có một phiếu. Vì vậy đa số là \textbf{No}.

\paragraph{Kết luận phân loại}
Với cách mã hoá và khoảng cách Euclidean đã nêu, k-NN với $K=5$ dự đoán cho khách mới: \textbf{không phê duyệt cho vay} (nhãn \textbf{No}).

\paragraph{Nhận xét ngắn (ổn định với $K$ nhỏ hơn)}
Nếu chọn $K=3$, ba láng giềng gần nhất vẫn là ID $5,8,7$ --- cả ba đều \textbf{No}, nên kết luận vẫn là \textbf{No}.

\paragraph{Chú ý độ đo \& chuẩn hoá (so sánh tùy chọn)}
Các chiều có độ lớn thang đo rất khác nhau (ví dụ số tiền vay nằm trên một trăm trong khi Gia đình chỉ có $0/1$), do đó nếu dùng thêm biến đổi Min--Max chỉ căn cứ vào các mẫu huấn luyện rồi lại đo Euclidean, với các $K$ trên một số láng giềng gần sẽ thay nhưng vẫn thu được đủ đa số \textbf{No}. Kết luận phân loại cuối cùng ở đây không phụ thuộc vào chuyển sang phương án đó trong bài làm tay này.
